\chapter{Hamiltonian actions on derived matrix spaces}
\label{ch:vol2-derived-hamiltonian-actions}

Imposing a matrix equation has two different effects on an action.
A bracket defect can vanish on its zero set, while the complex
that imposes the equation retains a homotopy for that defect.
The homotopies must themselves satisfy identities. For the
commuting-matrix equation, a calculation with cubic Hamiltonians
exhibits the first of these identities. A contraction of a free
associative resolution then solves them in every arity.

The resulting action preserves a cyclic orientation coherently.
Its matrix realization retains both the relation classes and the
gauge stabilizers. Requiring the homotopies to be Hamiltonian adds
a further condition: the scalar charges remember constants that
their vector fields forget. Quantizing the coefficient plane
introduces another distinction. Its finite-dimensional
representations and the quantization of matrix coordinates have
different equations, and consequently different recipients.

\section{The cubic defect and the relation complex}
\label{sec:vol2-matrix-cubic-defect}

Work over $\mathbb C$, with cohomological differentials of degree
one. All commutators are graded:
$[a,b]=ab-(-1)^{|a||b|}ba$ for homogeneous elements.
Put $A=\mathbb C[q,p]$ and give it the Poisson bracket
\[
 \{f,g\}=f_qg_p-f_pg_q.
\]
For a monomial $q^ap^b$, let $\sigma(q^ap^b)$ be the average of
the distinct words with $a$ occurrences of $q$ and $b$ occurrences of $p$.
Thus $\sigma(qp)=(qp+pq)/2$. The map $\sigma$ takes values in
the free associative algebra; it is a linear section of
abelianization.

Fix a positive integer $N$. Let $X,Y$ be $N\times N$ matrices of
independent commuting variables, with
\[
 \{X_{ij},Y_{kl}\}=\delta_{il}\delta_{jk},\qquad
 \{X_{ij},X_{kl}\}=\{Y_{ij},Y_{kl}\}=0.
\]
The scalar charge associated to $f\in A$ is
\begin{equation}\label{eq:vol2-matrix-charge}
 J_N(f)=\operatorname{Tr}\sigma f(X,Y).
\end{equation}
In particular $J_N(1)=N$. Write $\mu=[X,Y]=XY-YX$ for the
commuting-matrix equation. The brackets of cubic charges give
\begin{equation}\label{eq:vol2-cubic-trace-defect}
 \{J_N(q^3),J_N(p^3)\}-J_N(9q^2p^2)
   =-\frac32\operatorname{Tr}(\mu^2).
\end{equation}
Indeed, the first term is $9\operatorname{Tr}(X^2Y^2)$.
Of the six words in $\sigma(q^2p^2)$, four have cyclic trace
$\operatorname{Tr}(X^2Y^2)$ and two have cyclic trace
$\operatorname{Tr}(XYXY)$. Expanding $\operatorname{Tr}\mu^2$
proves the formula. At rank two, take
\[
 X=\begin{pmatrix}1&0\\0&0\end{pmatrix},\qquad
 Y=\begin{pmatrix}0&1\\1&0\end{pmatrix}.
\]
Then $\mu^2=-I_2$, so the defect is $3$.

The defect vanishes after imposing $\mu=0$. To retain how it
vanishes, introduce an odd matrix $\Psi$ of degree $-1$ and the
commutative differential graded algebra
\begin{equation}\label{eq:vol2-derived-matrix-algebra}
 R_N=\mathbb C[X_{ij},Y_{ij}]\otimes\Lambda(\Psi_{ij}),
 \qquad QX=QY=0,\qquad Q\Psi=[X,Y].
\end{equation}
Here and below, matrices with graded coefficients are multiplied
in their displayed order. The entries of $\Psi$ impose the
equations by a differential, including dependencies among them.
For example $\operatorname{Tr}\Psi$ is a cycle. It is not a
boundary: evaluation at $X=Y=0$ is a cochain map to
$\Lambda(\Psi_{ij})$ with zero differential and sends this cycle
to a nonzero element. At $N=1$ the differential vanishes entirely.
Thus the relation complex contains information beyond the
ordinary commuting scheme.

\section{A free resolution of the plane}
\label{sec:vol2-free-plane-resolution}

The homotopies can be constructed before choosing the matrix rank.
Consider the free differential graded associative algebra
\begin{equation}\label{eq:vol2-plane-resolution}
 R=\mathbb C\langle q,p,t\rangle,\qquad
 |q|=|p|=0,\quad |t|=-1,
 \qquad dq=dp=0,\quad dt=qp-pq.
\end{equation}
The projection $\pi:R\to A$ kills $t$ and makes $q,p$ commute.
Its linear section $i$ sends $q^ap^b$ to the normal word $p^bq^a$.
Both are cochain maps, but $i$ is not multiplicative:
$i(q)i(p)-i(qp)=dt$.

Replace the differential temporarily by $d_0t=qp$. Expand every
$t$ in a word into $qp$, marking the occurrence from which it
came. Two occurrences of $qp$ cannot overlap. For a fixed
expanded word, its marked subsets form a tensor product of
two-term complexes. The differential $d_0$ removes a mark, with
sign $(-1)^a$ if $a$ marks precede it.

Define $h_0$ by adding the first mark when that occurrence is
unmarked, and by zero when it is already marked or no occurrence
exists. Let $p_0$ send an unmarked normal word $p^bq^a$ to
$q^ap^b\in A$ and every other word to zero. The two-term
contraction gives
\begin{equation}\label{eq:vol2-first-word-contraction}
 d_0h_0+h_0d_0=1-ip_0,
 \qquad h_0^2=h_0i=p_0h_0=0.
\end{equation}
To see the signs, removing a later mark anticommutes with adding
the first. Removing the first mark contributes the identity.
When no occurrence exists, the word is already $p^bq^a$.

Write $d=d_0+\delta$, where $\delta t=-pq$. Exchanging an
expanded $qp$ for $pq$ decreases the number of pairs in which a
$q$ precedes a $p$. The operator $\delta h_0$ is therefore locally
nilpotent. Define
\begin{equation}\label{eq:vol2-word-contraction}
 h=h_0\sum_{j\geq0}(-\delta h_0)^j,
 \qquad p'=p_0\sum_{j\geq0}(-\delta h_0)^j.
\end{equation}
These are finite sums on every word.

\begin{proposition}\label{prop:vol2-plane-word-resolution}
The projection $p'$ equals $\pi$, and
\[
 \pi i=1,\qquad dh+hd=1-i\pi,
 \qquad h^2=hi=\pi h=0.
\]
Consequently $R\to A$ is a quasi-isomorphism and $h$ contracts
$J=\ker\pi$.
\end{proposition}
\begin{proof}
Substitute $(1+\delta h_0)^{-1}$ in
\eqref{eq:vol2-first-word-contraction}. Multiplying the finite
series by $1+\delta h_0$ cancels consecutive perturbation terms.
Since $\delta i=0$, the inclusion and the zero differential on
normal words remain unchanged. This gives
$dh+hd=1-ip'$ and the side identities.

The surviving reductions replace $qp$ by $pq$. A degree-zero
word therefore projects to its commutative monomial; a word
containing $t$ has negative degree and projects to zero. Thus
$p'=\pi$. On $J$ the homotopy identity becomes $dh+hd=1$.
\end{proof}

For a commutative differential graded algebra $B$, a morphism
$R\to\operatorname{Mat}_N(B)$ amounts to matrices $X,Y,\Psi$
with the differential in \eqref{eq:vol2-derived-matrix-algebra}.
This is the universal property represented by $R_N$.
The resolution of $A$ does not make $R_N$ acyclic in negative
degrees; the trace class already distinguishes the two assertions.

\section{The first action and its binary homotopy}
\label{sec:vol2-first-derived-action}

The Lie algebra of polynomial Hamiltonian vector fields is
$\mathfrak l=A/\mathbb C1$, with the induced Poisson bracket.
For a word $a_1\cdots a_m$ in $q,p$, define its cyclic derivative
by
\[
 \partial_q^{\mathrm{cyc}}(a_1\cdots a_m)
 =\sum_{a_j=q}a_{j+1}\cdots a_ma_1\cdots a_{j-1},
\]
and similarly for $p$. Removing an occurrence makes the remainder
start just after that occurrence. Cyclic rotation therefore does not
change the derivative.

For $f\in\mathfrak l$, define a degree-zero derivation of $R$ by
\begin{equation}\label{eq:vol2-hamiltonian-first-lift}
 F_1(f)q=-\partial_p^{\mathrm{cyc}}\sigma f,
 \qquad F_1(f)p=\partial_q^{\mathrm{cyc}}\sigma f,
 \qquad F_1(f)t=0.
\end{equation}
A constant has zero cyclic derivatives, so this is well defined.
The identity
\begin{equation}\label{eq:vol2-cyclic-derivative-identity}
 [q,\partial_q^{\mathrm{cyc}}W]
   +[p,\partial_p^{\mathrm{cyc}}W]=0
\end{equation}
follows by cancelling consecutive cyclic rotations of $W$.
It implies $[d,F_1(f)]=0$. Abelianization gives
$\pi F_1(f)=X_f\pi$, where $X_f=\{f,-\}$.

Set $\mu_R=qp-pq$. Cyclic differentiation gives
\begin{equation}\label{eq:vol2-cubic-vector-defect}
 \bigl([F_1(q^3),F_1(p^3)]-F_1(9q^2p^2)\bigr)(q,p,t)
   =\bigl(3[\mu_R,q],3[\mu_R,p],0\bigr).
\end{equation}
The derivation
\begin{equation}\label{eq:vol2-cubic-word-primitive}
 F_2(q^3,p^3)(q,p,t)
     =\bigl(3[t,q],3[t,p],3t^2\bigr)
\end{equation}
has degree $-1$ and has this defect as its differential
$[d,F_2]$. On $q,p$ the assertion follows from $dt=\mu_R$.
On $t$, the identity
\[
 d(t^2)=\mu_Rt-t\mu_R
\]
cancels the term $F_2(\mu_R)$. The value on $t$ is therefore
necessary even though the original defect vanishes on $t$.

The failure of strictness cannot be removed by another polynomial
choice of cyclic words. Let
\[
 \mathcal N=
 \mathbb C\langle q,p\rangle/
 \bigl(\mathbb C1+[\mathbb C\langle q,p\rangle,
                         \mathbb C\langle q,p\rangle]\bigr),
\]
where the commutators span a vector subspace. Its elements are
cyclic words modulo constants. The necklace bracket is
\[
 \{[U],[V]\}_{\mathcal N}
   =\bigl[\partial_q^{\mathrm{cyc}}U\,
               \partial_p^{\mathrm{cyc}}V
          -\partial_p^{\mathrm{cyc}}U\,
               \partial_q^{\mathrm{cyc}}V\bigr].
\]
This bracket has a useful derivation description. Send a cyclic
word $U$ to $D_Uq=-\partial_p^{\mathrm{cyc}}U$ and
$D_Up=\partial_q^{\mathrm{cyc}}U$. Differentiating the displayed
bracket gives
$[D_U,D_V]=D_{\{[U],[V]\}_{\mathcal N}}$: terms selecting
two positions in one cyclic word cancel in pairs, while terms
selecting one position in each word give the two derivation
composites. This also proves that the bracket is well defined
under cyclic rotation. The map to derivations is injective
modulo constants, since for a homogeneous word of length $r>0$,
\[
 [q\,\partial_q^{\mathrm{cyc}}U
   +p\,\partial_p^{\mathrm{cyc}}U]=r[U].
\]
Jacobi follows from Jacobi for derivations and this injectivity.
Abelianization is a Lie map $\mathcal N\to\mathfrak l$.

\begin{proposition}\label{prop:vol2-necklace-no-section}
The abelianization $\mathcal N\to\mathfrak l$ has no linear
Lie section, including sections that do not preserve polynomial
degree.
\end{proposition}
\begin{proof}
In degrees at most three there is only one cyclic word for each
commutative monomial. A lift of a cubic therefore differs from
its cyclic cubic by terms of degree at least four. Bracketing
such a correction with a cubic has degree at least five.

Put $C=[q^2p^2]$ and $D=[qpqp]$. They are distinct cyclic words,
and direct differentiation gives
\[
 \{[q^3],[p^3]\}_{\mathcal N}=9C,
 \qquad \{[q^2p],[qp^2]\}_{\mathcal N}=C+2D.
\]
The commutative brackets are respectively $9q^2p^2$ and
$3q^2p^2$. The degree-four part of a Lie section applied to
$q^2p^2$ would have to equal both $C$ and $(C+2D)/3$.
Higher-degree corrections cannot change this contradiction.
\end{proof}

This obstruction concerns a strict polynomial necklace section.
The relation generator in \eqref{eq:vol2-plane-resolution}
provides the larger differential graded carrier on which
\eqref{eq:vol2-cubic-word-primitive} solves the first defect.

\section{Solving the action equations in every arity}
\label{sec:vol2-all-arity-derived-action}

Give $\operatorname{Der}(R)$ its graded commutator and differential
$\partial D=[d,D]$. The derivations with zero projection form
a differential graded Lie algebra
\[
 K=\{D:\pi D=0\}.
\]
A degree-$j$ derivation in $K$ is specified by
$U=Dq,V=Dp\in J^j$ and $W=Dt\in J^{j-1}$.
Put $L(U,V)=[U,p]+[q,V]$. Then
\begin{equation}\label{eq:vol2-derivation-differential}
 (\partial D)(q,p,t)
   =\bigl(dU,dV,dW-(-1)^jL(U,V)\bigr).
\end{equation}
The contraction of words yields a contraction of this entire
derivation complex, including its nonclosed elements.

\begin{lemma}\label{lem:vol2-derivation-kernel-contraction}
The degree-minus-one map
\begin{equation}\label{eq:vol2-derivation-kernel-homotopy}
 (HD)(q,p,t)
   =\bigl(hU,hV,h(W-(-1)^jL(hU,hV))\bigr)
\end{equation}
preserves $K$ and satisfies $\partial H+H\partial=1_K$.
Every value of $H$ is a finite polynomial.
\end{lemma}
\begin{proof}
The ideal $J$ contains all arguments of $h$ in the formula.
Since $\pi h=0$, the resulting derivation again lies in $K$.
On $q,p$ the contraction identity is $dh+hd=1_J$.
Use $dL(U,V)=L(dU,dV)$ to substitute the formula on $t$.
The sum of the third components is
\[
 (dh+hd)W+
 (-1)^j hL\bigl((dh+hd-1)U,(dh+hd-1)V\bigr)=W.
\]
The word contraction is finite, so the formula has the stated
polynomial meaning.
\end{proof}

We specify all higher signs by a coalgebra convention.
For a differential graded Lie algebra $\mathfrak g$, let $s$
be the degree-minus-one shift to $\mathfrak g[1]$.
The cofree conilpotent symmetric coalgebra
$S^c(\mathfrak g[1])=\bigoplus_{n\geq0}S^n(\mathfrak g[1])$
has the unshuffle coproduct. Its Lie codifferential has
corestrictions
\begin{equation}\label{eq:vol2-lie-coalgebra-signs}
 b_1(sa)=-s\partial a,
 \qquad b_2(sa\,sb)=(-1)^{|a|}s[a,b].
\end{equation}
An $L_\infty$ morphism is a coaugmented coalgebra map commuting
with these codifferentials. On degree-zero source inputs, its
arity-$n$ corestriction is $sF_n$, where
$F_n:\Lambda^n\mathfrak l\to\operatorname{Der}(R)$ has degree
$1-n$.

Suppose $F_1,\ldots,F_{n-1}$ are defined. Their corestrictions
extend by set partitions to a coalgebra map $\mathcal F$ with
arity-$n$ corestriction zero. Define $\Omega_n$ by
\[
 \operatorname{pr}_1(b\mathcal F-\mathcal F b)\big|_{S^n}
     =s\Omega_n,
 \qquad F_n=H\Omega_n.
\]
For the first two residuals this says
\begin{align}
 \Omega_2(f,g)
    &=[F_1(f),F_1(g)]-F_1(\{f,g\}),
       \label{eq:vol2-action-second-residual}\\
 \Omega_3(f,g,k)
    &=\sum_{\mathrm{cyc}}
       \bigl([F_1(f),F_2(g,k)]-F_2(\{f,g\},k)\bigr).
       \label{eq:vol2-action-third-residual}
\end{align}
At arity four, the brackets between two binary components contribute
\begin{equation}\label{eq:vol2-action-fourth-binary-part}
 \begin{aligned}
 &-[F_2(f_1,f_2),F_2(f_3,f_4)]
   +[F_2(f_1,f_3),F_2(f_2,f_4)]\\
 &\hspace{3em}-[F_2(f_1,f_4),F_2(f_2,f_3)].
 \end{aligned}
\end{equation}
They accompany the unary--ternary and source-bracket terms.

\begin{theorem}\label{thm:vol2-all-arity-derived-action}
The specified first action and the recursion $F_n=H\Omega_n$
define an $L_\infty$ morphism
\[
 F:\mathfrak l\longrightarrow\operatorname{Der}(R).
\]
Every component has finite polynomial values on polynomial
inputs. For every $N\geq1$, matrix substitution gives a
$GL_N$-equivariant $L_\infty$ action on $R_N$, and hence on
$R_N^{GL_N}$.
\end{theorem}
\begin{proof}
The first component commutes with $d$ by
\eqref{eq:vol2-cyclic-derivative-identity}. Let
$E=b\mathcal F-\mathcal F b$. The identity
$bE+Eb=0$ follows from the square-zero codifferentials.
If $E$ vanishes below arity $n$, its arity-$n$ corestriction
therefore satisfies $\partial\Omega_n=0$.

The first action projects to the strict Hamiltonian action on
$A$, so $\Omega_2$ projects to zero. Moreover $F_1$ normalizes
$K$, and every constructed component of arity at least two
belongs to $K$. The same is consequently true of each subsequent
residual. Lemma~\ref{lem:vol2-derivation-kernel-contraction}
gives $\partial H\Omega_n=\Omega_n$. In the coalgebra equation,
the new term is $-s\partial F_n$, by
\eqref{eq:vol2-lie-coalgebra-signs}; it cancels $s\Omega_n$.
Induction proves all the identities.

At each arity there are finitely many partitions. The word
contraction terminates on every term, proving polynomial
finiteness. A word derivation acts on matrix entries by
substituting $X,Y,\Psi$ for $q,p,t$ in its generator values.
The graded Leibniz rule makes this substitution preserve
differentials and brackets. It also commutes with simultaneous
conjugation, which proves the two matrix assertions.
\end{proof}

The contraction gives explicit higher terms beyond the binary
homotopy. For
$\mathbf f=(q^3,p^3,qp^3)$ and
$\mathbf g=(q^3,p^3,qp^3,q^2p^2)$, one obtains
\begin{equation}\label{eq:vol2-derived-higher-examples}
 \begin{aligned}
 F_3(\mathbf f)q&=-\frac92(tpt+pt^2),\\
 F_4(\mathbf g)q&=\frac{33}{5}t^3p+\frac{39}{5}t^2pt
                   +\frac{69}{5}tpt^2+\frac{99}{5}pt^3.
 \end{aligned}
\end{equation}
The other generator values are
$F_3(\mathbf f)p=F_3(\mathbf f)t=0$ and
$F_4(\mathbf g)p=F_4(\mathbf g)t=0$.
These formulas follow by applying
\eqref{eq:vol2-derivation-kernel-homotopy} to the residuals.
They can also be checked by substituting $dt=qp-pq$ in
$\partial F_n=\Omega_n$ on each generator. The four-input
identity uses \eqref{eq:vol2-action-fourth-binary-part}.
Thus this action does not truncate at its binary term.
The displayed values depend on the contraction; a nonzero
chosen component alone does not establish an invariant
obstruction to changing the action's presentation.

\section{Choices and pointed formal Hamiltonians}
\label{sec:vol2-derived-action-choices}

Let $\operatorname{Der}_\pi(R)$ be the derivations preserving
$J=\ker\pi$. Projection gives an exact sequence of differential
graded Lie algebras
\begin{equation}\label{eq:vol2-derived-action-fibration}
 0\longrightarrow K\longrightarrow\operatorname{Der}_\pi(R)
   \xrightarrow{\pi_*}\operatorname{Der}(A)\longrightarrow0.
\end{equation}
Surjectivity is degreewise. A derivation of $A$, concentrated in
degree zero, lifts by choosing words for its values on $q,p$
and giving value zero on $t$. Such a lift preserves $J$, although
it need not commute with $d$. The contracted kernel shows that
$\pi_*$ is a quasi-isomorphism.

For a differential graded Lie algebra $L$, retain homotopies
between actions by using the arity-completed
convolution differential graded Lie algebra
\[
 \mathcal G(L)=\prod_{n\geq1}
    \operatorname{Hom}(S^n(\mathfrak l[1]),L),
 \qquad
 \mathcal F^r\mathcal G(L)=\prod_{n\geq r}
    \operatorname{Hom}(S^n(\mathfrak l[1]),L).
\]
The differential combines the source Chevalley codifferential
with that of $L$; the bracket uses the reduced coproduct.
For a degree-one element $\alpha$, the Maurer--Cartan equation
$\partial\alpha+[\alpha,\alpha]/2=0$ is the $L_\infty$
morphism equation in the convention
\eqref{eq:vol2-lie-coalgebra-signs}. Polynomial differential
forms on a simplex give parameter families of such equations.
Tensoring arity by arity defines their simplicial space.

The \emph{derived space of lifts} is the homotopy fibre, over
$f\mapsto X_f$, of the map of these spaces induced by $\pi_*$.
This fixes the action on $A$ and the comparison to it; it does
not yet impose a cyclic or Hamiltonian condition on the lifts.

\begin{proposition}\label{prop:vol2-hamiltonian-lift-space}
The derived space of lifts of $X$ to
$\operatorname{Der}_\pi(R)$ is nonempty and contractible.
The word contraction selects one point of this space.
\end{proposition}
\begin{proof}
On the arity-$n$ associated graded, the kernel is
\[
 \operatorname{Hom}(S^n(\mathfrak l[1]),K),
\]
contracted by postcomposition with $H$, with the tensor sign
when parameter forms are present. Polynomial forms on a simplex
restrict degreewise onto compatible forms on a boundary or horn.
One proves this by extending faces successively: coefficients
vanishing on earlier faces lie in the corresponding ideals
$(x_i,dx_i)$, and the barycentric relation $\sum x_i=1$
handles the last face. No cochain primitive in the parameter
direction is required.

Suppose the lower arities and a compatible boundary family are
already lifted. Choose an arity-$n$ component with that boundary
and the prescribed projection. Its curvature lies in the kernel
above, tensored with forms relative to the boundary. The
arity-$n$ Bianchi identity makes the curvature closed. The
contraction gives a relative primitive; subtracting it solves
the equation without changing the boundary or projection.
The same argument fills horns. Induction and arity completeness
give a trivial fibration of simplicial spaces. Its nonempty
fibre is contractible.
\end{proof}

There is also a completion in polynomial weight. Give $q,p$
weight one and $t$ weight two; write $R_w$ for the homogeneous
piece and put $\widehat R=\prod_{w\geq0}R_w$.
The product topology has discrete homogeneous pieces. Assign
the same weights to $X,Y,\Psi$ to define $\widehat R_N$.
A pointed Hamiltonian has no constant or linear term.

\begin{proposition}\label{prop:vol2-pointed-action-completion}
On homogeneous inputs of degrees $d_1,\ldots,d_n$, the
derivation $F_n$ shifts weight by
\[
 \sum_{i=1}^n(d_i-2).
\]
At each arity, the action therefore extends continuously from
pointed polynomial Hamiltonians to
\[
 \widehat{\mathfrak l}_{\mathrm{pt}}
   =\prod_{d\geq2}\operatorname{Sym}^d(\mathbb Cq\oplus\mathbb Cp)
\]
on $\widehat R$, and to the corresponding equivariant action
on $\widehat R_N$.
\end{proposition}
\begin{proof}
Both $d$ and $h$ preserve weight. Cyclic differentiation gives
the claimed shift for $F_1$. Brackets add shifts, and the
derivation contraction preserves them, so induction proves the
formula for every component.

For pointed inputs, each $d_i-2$ is nonnegative. At a fixed
arity and output weight, only finitely many source weights
and input degrees can contribute. Their homogeneous spaces
are finite dimensional. The operation and each of its identities
are thus finite in every output weight, giving the continuous
extension. Matrix substitution preserves these weights.
\end{proof}

\section{Gauge ghosts and Hamiltonian primitives}
\label{sec:vol2-derived-action-bfv}

The matrices also carry simultaneous conjugation. Introduce
an odd matrix $c$ of degree $1$ and the polynomial graded
commutative algebra
\[
 \mathcal B_N=\mathbb C[X_{ij},Y_{ij}]
                  \otimes\Lambda(c_{ij},\Psi_{ij}).
\]
Extend the even Poisson bracket by
\[
 \{c_{ij},\Psi_{kl}\}=\{\Psi_{kl},c_{ij}\}
    =\delta_{il}\delta_{jk},
\]
with all other new generator brackets zero. The symmetry on
these two odd generators is the graded skew-symmetry of an
even Poisson bracket. The degree-one Hamiltonian
\[
 \Theta=\operatorname{Tr}(c[X,Y]-\Psi c^2)
\]
defines $Q=\{\Theta,-\}$, with
\begin{equation}\label{eq:vol2-classical-bfv-differential}
 QX=-[c,X],\quad QY=-[c,Y],\quad Qc=-c^2,
 \quad Q\Psi=[X,Y]-[c,\Psi].
\end{equation}
Here $[c,\Psi]=c\Psi+\Psi c$. Applying the graded Leibniz
rule once more to the four expressions gives $Q^2=0$.
This is the classical BFV, or Hamiltonian BRST, complex.
Equivalently it is the Chevalley--Eilenberg algebra
$C^*_{\mathrm{CE}}(\mathfrak{gl}_N;R_N)$ with the displayed
gauge sign. Its positive ghosts retain the infinitesimal gauge
action in addition to the negative relation generators.

\begin{proposition}\label{prop:vol2-action-bfv-extension}
Every component of the matrix action extends to $\mathcal B_N$
by giving value zero on $c$. These extensions define an
$L_\infty$ action of $\mathfrak l$ on $\mathcal B_N$.
\end{proposition}
\begin{proof}
Let $D$ be a word derivation of degree $j$ with $Dc=0$.
For a homogeneous matrix word $M$, the graded Leibniz rule gives
$D[c,M]=(-1)^j[c,DM]$. Thus the extended derivation commutes
in the graded sense with the gauge part of
\eqref{eq:vol2-classical-bfv-differential}. Extension preserves
both the differential and the bracket of derivations, so it
preserves every action identity.
\end{proof}

The first derivation is Hamiltonian: it equals
$H_{J_N(f)}=\{J_N(f),-\}$. Each $J_N(f)$ is $Q$-closed by
conjugation invariance. The cubic scalar defect has a
Hamiltonian primitive
\begin{equation}\label{eq:vol2-cubic-hamiltonian-primitive}
 K_N=-\frac32\operatorname{Tr}(\Psi\mu),\qquad
 QK_N=-\frac32\operatorname{Tr}(\mu^2).
\end{equation}
Graded cyclicity cancels the gauge terms in $QK_N$;
\eqref{eq:vol2-cubic-trace-defect} identifies its remaining
term with the bracket defect. Its Hamiltonian vector field is
\begin{equation}\label{eq:vol2-cubic-hamiltonian-field}
 \begin{aligned}
 H_{K_N}X&=\tfrac32[\Psi,X],&
 H_{K_N}Y&=\tfrac32[\Psi,Y],\\
 H_{K_N}c&=-\tfrac32\mu,&H_{K_N}\Psi&=0.
 \end{aligned}
\end{equation}
The extended word homotopy, in contrast, has
\begin{equation}\label{eq:vol2-cubic-derived-field}
 DX=3[\Psi,X],\quad DY=3[\Psi,Y],
 \quad Dc=0,\quad D\Psi=3\Psi^2.
\end{equation}
Both satisfy
\[
 [Q,D]=[Q,H_{K_N}]
   =[H_{J_N(q^3)},H_{J_N(p^3)}]-H_{J_N(9q^2p^2)}.
\]
The ghost component in \eqref{eq:vol2-cubic-hamiltonian-field}
is part of its Hamiltonian solution.

For $N\geq2$, $D$ is not Hamiltonian for the displayed bracket.
Indeed, $Dc=0$ would force a polynomial Hamiltonian for $D$ to
be independent of every $\Psi_{ij}$. Such a polynomial has no
component of degree $-1$, while $D$ is a nonzero degree-minus-one
derivation. A coherent derivation action therefore does not by
itself supply Hamiltonian representatives for its homotopies.

There is a prior scalar obstruction to a homotopy moment map
on the quotient $\mathfrak l$. Choose representatives with
constant term zero. Their bracket can acquire a constant:
the scalar cocycle is $\kappa(f,g)=\{f,g\}(0)$.

\begin{proposition}\label{prop:vol2-moment-central-obstruction}
For $N>0$, no $L_\infty$ morphism
$\mathfrak l\to(\mathcal B_N,Q,\{-,-\})$ has first component
$f\mapsto J_N(f)$ on those representatives. The obstruction
persists if the first components are changed by $Q$-boundaries.
\end{proposition}
\begin{proof}
The classes of $q,p$ have zero bracket in $\mathfrak l$, whereas
\[
 \{\operatorname{Tr}X,\operatorname{Tr}Y\}=N.
\]
The binary morphism equation would make $N$ a $Q$-boundary.
Evaluation of $X,Y,c,\Psi$ at zero is a cochain map to
$\mathbb C$, since every value in
\eqref{eq:vol2-classical-bfv-differential} is quadratic.
It sends the alleged boundary to the nonzero scalar $N$.
Exact changes of closed charges leave their cohomology bracket
unchanged, so they cannot remove this obstruction.
\end{proof}

Retaining the central extension $A\to\mathfrak l$, with
$J_N(1)=N$, removes this particular obstruction. The cocycle
also vanishes on the classical pointed subalgebra $(q,p)^2$.
On vector fields constants already act by zero, which explains
why Theorem~\ref{thm:vol2-all-arity-derived-action} is consistent
with the proposition.

The central extension contains a strict finite sector. Write
$e=q^2/2$, $f=p^2/2$ and $h_{\mathrm{quad}}=qp$.
Besides skew versions, its
nonzero brackets are
\[
 \{q,p\}=1,\quad\{e,p\}=q,\quad\{f,q\}=-p,
 \quad\{e,f\}=h_{\mathrm{quad}},
\]
\[
 \{h_{\mathrm{quad}},q\}=-q,\quad
 \{h_{\mathrm{quad}},p\}=p,\quad
 \{h_{\mathrm{quad}},e\}=-2e,\quad
 \{h_{\mathrm{quad}},f\}=2f.
\]
The map $J_N$ sends $1,q,p,e,h_{\mathrm{quad}},f$ to
$N,\operatorname{Tr}X,\operatorname{Tr}Y,
\operatorname{Tr}(X^2)/2,\operatorname{Tr}(XY),
\operatorname{Tr}(Y^2)/2$. Canonical matrix differentiation
gives precisely the same list of brackets. Thus scalar
quadratics act by a strict Hamiltonian Lie map of closed
functions. Higher scalar Hamiltonians require the homotopies
whose first instance is \eqref{eq:vol2-cubic-hamiltonian-primitive}.

\section{A cyclic orientation of the resolution}
\label{sec:vol2-derived-cyclic-orientation}

The symplectic plane has the volume form $dq\wedge dp$.
To preserve it on a resolution, the action must act on a complex
that remembers cyclic order as well as the differential.
Use normalized Hochschild chains
\[
 C_n(R)=R\otimes\overline R^{\otimes n},\qquad
 \overline R=R/\mathbb C1,
\]
written $a_0[a_1|\cdots|a_n]$. A unit in any bar slot represents
zero. A chain of internal degree $m=\sum_i|a_i|$ has total
cohomological degree $m-n$. Let $d_\otimes$ be the tensor
differential with its ordinary Koszul signs. In the raw-chain
convention the Hochschild differential is
\begin{equation}\label{eq:vol2-raw-hochschild-differential}
 d_{\mathrm H}=(-1)^n d_\otimes+b,
\end{equation}
where
\begin{align*}
 b(a_0[a_1|\cdots|a_n])
  ={}&a_0a_1[a_2|\cdots|a_n]\\
    &+\sum_{j=1}^{n-1}(-1)^j
        a_0[a_1|\cdots|a_ja_{j+1}|\cdots|a_n]\\
    &+(-1)^{n+|a_n|\sum_{i<n}|a_i|}
        a_na_0[a_1|\cdots|a_{n-1}].
\end{align*}
For $n=0$ set $b=0$. The normalized Connes operator is
\begin{equation}\label{eq:vol2-raw-connes-operator}
 B(a_0[a_1|\cdots|a_n])
 =\sum_{j=0}^n
 (-1)^{nj+(\sum_{i<j}|a_i|)(\sum_{i\geq j}|a_i|)}
 1[a_j|\cdots|a_n|a_0|\cdots|a_{j-1}].
\end{equation}
Adjacent face cancellations and cyclic rotation give
$d_{\mathrm H}^2=B^2=d_{\mathrm H}B+Bd_{\mathrm H}=0$.
Thus $(C_\bullet(R),d_{\mathrm H},B)$ is a mixed complex,
with the second operator of degree $-1$.

Let $\zeta$ have degree two. The negative-cyclic complex is
\[
 M_R=\bigl(C_\bullet(R)[[\zeta]],d_{\mathrm H}+\zeta B\bigr).
\]
Its degree-minus-two orientation is
\begin{equation}\label{eq:vol2-raw-orientation-cycle}
 \omega_0=1[q|p]-1[p|q]-1[t].
\end{equation}
The first two terms have Hochschild boundary $-1[qp-pq]$.
The internal differential of the last term, with the sign in
\eqref{eq:vol2-raw-hochschild-differential}, is $+1[qp-pq]$.
Its remaining Hochschild boundary is zero. Every Connes term
has a unit in a bar slot. Hence
$d_{\mathrm H}\omega_0=B\omega_0=0$.

The sign of $1[t]$ is tied to this totalization. Multiplication
of a chain of bidegree $(m,n)$ by $(-1)^{mn}$ changes the
differential to $d_\otimes+(-1)^m b$ and the Connes operator
to $(-1)^mB$; in that convention the last sign is positive.
The explicit formulas below use the raw convention throughout.

Let $\Omega_A^a$ denote polynomial differential $a$-forms with
de Rham differential $d_{\mathrm{dR}}$, and place it in
cohomological degree $-a$. The normalized
Hochschild--Kostant--Rosenberg projection is
\begin{equation}\label{eq:vol2-plane-hkr-projection}
 p_{\mathrm H}(a_0[a_1|\cdots|a_n])
  =\frac1{n!}\pi(a_0)\,d_{\mathrm{dR}}\pi(a_1)
                     \wedge\cdots\wedge d_{\mathrm{dR}}\pi(a_n).
\end{equation}
Substitution in the face and cyclic formulas gives
\[
 p_{\mathrm H}d_{\mathrm H}=0,
 \qquad p_{\mathrm H}B=d_{\mathrm{dR}}p_{\mathrm H},
 \qquad p_{\mathrm H}\omega_0=dq\wedge dp.
\]
The factorial is essential to the last normalization.

\begin{proposition}\label{prop:vol2-mixed-plane-contraction}
There is a weight-preserving contraction
$(p_{\mathrm H},i_{\mathrm H},H_{\mathrm H})$ from
$(C_\bullet(R),d_{\mathrm H})$ to $(\Omega_A^\bullet,0)$,
with $i_{\mathrm H}(dq\wedge dp)=\omega_0$ and
\[
 H_{\mathrm H}^2=H_{\mathrm H}i_{\mathrm H}
                =p_{\mathrm H}H_{\mathrm H}=0.
\]
Its negative-cyclic perturbation has unchanged projection and
target differential $\zeta d_{\mathrm{dR}}$. More precisely,
\begin{equation}\label{eq:vol2-mixed-plane-perturbation}
 \begin{aligned}
 i_\zeta&=(1+\zeta H_{\mathrm H}B)^{-1}i_{\mathrm H},\\
 H_\zeta&=H_{\mathrm H}(1+\zeta BH_{\mathrm H})^{-1},
       &p_\zeta&=p_{\mathrm H}.
 \end{aligned}
\end{equation}
The choices can be made by finite rational linear algebra at
each polynomial weight.
\end{proposition}
\begin{proof}
Give $q,p,t$ weights $1,1,2$ and each bar chain the sum of its
generator weights. Normalized bar slots have positive weight, so
the bar-length filtration is finite in each fixed weight.
The quasi-isomorphism $R\to A$ therefore gives a
quasi-isomorphism of Hochschild chains weight by weight:
on each associated-graded piece the comparison is a tensor
product of quasi-isomorphisms over the field.

Write $A^e=A\otimes A^{\mathrm{op}}$, and use subscripts $L,R$
for its two factors. The diagonal of $A$ has the bimodule
Koszul resolution
\[
 A^e\otimes\Lambda(e_q,e_p),\qquad
 de_q=q_L-q_R,\qquad de_p=p_L-p_R.
\]
The two differences form a regular sequence, since they become
independent polynomial variables after a linear change of
coordinates. Tensoring with $A$ over $A^e$ kills this
differential and identifies $e_q,e_p$ with $dq,dp$.
The bar resolution is another free resolution of the diagonal;
its augmented contraction inserts the unit. The comparison sends
$e_q\wedge e_p$ to the antisymmetrized bar
$1[q|p]-1[p|q]$. Consequently
\eqref{eq:vol2-plane-hkr-projection} is a surjective
quasi-isomorphism with the stated normalization.

Choose lifts of a basis of forms to cycles, using $\omega_0$
for the volume. These give $i_{\mathrm H}$. In the acyclic
kernel, choose degreewise splittings
\[
 (\ker p_{\mathrm H})^r=B^r\oplus L^r,
 \qquad d_{\mathrm H}:L^r\xrightarrow{\sim}B^{r+1}.
\]
Take $H_{\mathrm H}$ to be the inverse on $B^{r+1}$ and zero
on $L$ and the chosen cycle lifts. This proves the contraction
and side identities. Every space in a fixed weight is finite
dimensional; lexicographic word bases and row reduction make
the choices rational.

Apply perturbation by $\zeta B$. Its series converge in the
$\zeta$-adic topology. The relation
\[
 p_{\mathrm H}\zeta BH_{\mathrm H}
      =\zeta d_{\mathrm{dR}}p_{\mathrm H}H_{\mathrm H}=0
\]
keeps the projection fixed. It also kills every higher term
in the transferred differential
$p_{\mathrm H}\zeta B(1+\zeta H_{\mathrm H}B)^{-1}i_{\mathrm H}$,
leaving exactly $\zeta d_{\mathrm{dR}}$.
Since $B\omega_0=0$, the perturbed inclusion still sends the
volume to $\omega_0$.
\end{proof}

There is no unbounded filtration exchange in this comparison.
At fixed weight the bar length and the negative internal degree
are bounded. Fixing total degree then bounds the exponent of
$\zeta$. The completed identities can therefore be checked in
finite-dimensional pieces as well as in finite parameter quotients.

\section{Coherent preservation of the volume}
\label{sec:vol2-equivariant-cyclic-volume}

A derivation $D$ of degree $r$ acts on raw chains by
\begin{equation}\label{eq:vol2-raw-chain-lie-derivative}
 \mathcal L_D(a_0[a_1|\cdots|a_n])
  =\sum_{j=0}^n(-1)^{nr+r\sum_{i<j}|a_i|}
       a_0[\cdots|D(a_j)|\cdots|a_n],
\end{equation}
where the $j=0$ term acts on the outer coefficient. The factor
$(-1)^{nr}$ is required by the raw totalization. The tensor
signs give
\[
 [d_{\mathrm H},\mathcal L_D]=\mathcal L_{[d,D]},
 \qquad [B,\mathcal L_D]=0,
 \qquad [\mathcal L_D,\mathcal L_E]=\mathcal L_{[D,E]}.
\]
Thus the action $F$ makes $M_R$ an $L_\infty$ module.

For explicit totalization, write $C=S^c(\mathfrak l[1])$ and
let $b_C$ be its Chevalley codifferential. The twisting cochain
$\alpha:C\to\operatorname{Der}(R)$ has degree one and
\[
 \alpha_n(sf_1\cdots sf_n)=F_n(f_1,\ldots,f_n).
\]
Its equation is
$\partial\alpha+\alpha b_C+[\alpha,\alpha]/2=0$,
where convolution uses the tensor Koszul sign. On the completed
coefficient complex
\[
 \mathcal C_R=\prod_{n\geq0}
       \operatorname{Hom}(S^n(\mathfrak l[1]),M_R),
\]
the differential of a homogeneous map $a$ is
\begin{equation}\label{eq:vol2-equivariant-cyclic-differential}
 D_Fa=d_{M_R}a-(-1)^{|a|}a b_C+\mathcal L_\alpha\star a.
\end{equation}
Here $\star$ denotes action convolution with the unshuffle
coproduct.
The twisting-cochain equation implies $D_F^2=0$.

The projection to forms intertwines the first action and kills
the higher ones:
\begin{equation}\label{eq:vol2-equivariant-hkr}
 p_{\mathrm H}\mathcal L_{F_1(f)}
      =\mathcal L_{X_f}p_{\mathrm H},
 \qquad p_{\mathrm H}\mathcal L_{F_n}=0\quad(n\geq2).
\end{equation}
The first equation is naturality of the differential-form
formula. For the second, a negative-degree value of a
derivation contains $t$ and is killed by $\pi$.

A coherent lift is a $D_F$-closed element of total degree
$-2$ whose projection is the volume, regarded as a constant
arity-zero cochain. Its arity-$n$ component is an alternating
map $\omega_n:\Lambda^n\mathfrak l\to M_R^{-n-2}$.
Polynomial parameter forms retain homotopies over this fixed
projection and define the derived space of such lifts.

\begin{theorem}\label{thm:vol2-coherent-cyclic-orientation}
The orientation $\omega_0$ extends to a $D_F$-closed element
$\omega_{\mathrm{eq}}$ of total degree $-2$.
Over the fixed Hamiltonian action and fixed projection to
$dq\wedge dp$, the derived space of coherent negative-cyclic
lifts is contractible. The construction is continuous for
pointed Hamiltonians in the positive-weight completion.
\end{theorem}
\begin{proof}
Extend the contraction of
Proposition~\ref{prop:vol2-mixed-plane-contraction} to the
convolution complex, with its total signs. Let
$\mathcal H_\zeta$ denote the extended homotopy and let
$\delta_F$ denote the terms in
\eqref{eq:vol2-equivariant-cyclic-differential} beyond the
coefficient differential $d_{M_R}$. They raise arity.
Perturbation in the arity filtration gives
\begin{equation}\label{eq:vol2-equivariant-volume-formula}
 \omega_{\mathrm{eq}}
    =(1+\mathcal H_\zeta\delta_F)^{-1}
          i_\zeta(dq\wedge dp).
\end{equation}
Equation~\eqref{eq:vol2-equivariant-hkr} and
$p_{\mathrm H}\mathcal H_\zeta=0$ imply
$p_{\mathrm H}\delta_F\mathcal H_\zeta=0$.
The transferred differential is therefore precisely the strict
Hamiltonian action on forms together with
$\zeta d_{\mathrm{dR}}$. The volume is closed and its Lie
derivative under $X_f$ is zero, since
$\partial_q(-f_p)+\partial_p(f_q)=0$.
Its perturbed inclusion is consequently closed. Its arity-zero
component is $\omega_0$.

The same perturbation contracts the kernel of the fixed
projection. The homotopy fibre over the specified closed volume
is an affine translate of this acyclic complex. Polynomial
parameter forms and its contraction fill relative families,
so its derived mapping space is contractible.

At fixed arity only finitely many arity-raising factors occur.
At fixed $\zeta$ precision the mixed perturbation is finite.
Finally, the weight estimate in
Proposition~\ref{prop:vol2-pointed-action-completion} makes
every output weight a finite sum for pointed inputs. These
three finiteness statements prove the completion assertion.
\end{proof}

Two coherence levels can be written without a matrix inversion.
Abbreviate $1[a_1|\cdots|a_n]$ by
$\langle a_1|\cdots|a_n\rangle$. The arity-one equation is
\[
 d_{\mathrm H}\omega_1(f)=-\mathcal L_{F_1(f)}\omega_0.
\]
Choose
\begin{equation}\label{eq:vol2-cyclic-cubic-primitives}
 \omega_1(q^3)=-3\langle q|q|q\rangle,
 \qquad \omega_1(p^3)=-3\langle p|p|p\rangle.
\end{equation}
For instance,
$d_{\mathrm H}\langle q|q|q\rangle
=\langle q|q^2\rangle-\langle q^2|q\rangle$.
A compatible quartic prescription is
\begin{align}
 \omega_1(q^2p^2)
 ={}&-\tfrac23\langle tt\rangle
       \label{eq:vol2-cyclic-quartic-primitive}\\
 &+\tfrac23\bigl(
       \langle q|pt\rangle-\langle q|tp\rangle
      -\langle p|qt\rangle+\langle p|tq\rangle\bigr)\notag\\
 &+\tfrac23\bigl(
       \langle qt|p\rangle-\langle pt|q\rangle
      -\langle tq|p\rangle+\langle tp|q\rangle\bigr)\notag\\
 &-2\bigl(\langle q|p|qp\rangle
           +\langle p|qp|q\rangle+\langle qp|q|p\rangle\bigr).
       \notag
\end{align}
All three chains have zero Connes boundary and zero HKR
projection. For $f=q^3,g=p^3$, the second equation in the same
convention is
\begin{align}
 d_{\mathrm H}\omega_2(f,g)
 ={}&\omega_1(\{f,g\})
       -\mathcal L_{F_1(f)}\omega_1(g)
       +\mathcal L_{F_1(g)}\omega_1(f)\notag\\
    &-\mathcal L_{F_2(f,g)}\omega_0.
       \label{eq:vol2-cyclic-second-equation}
\end{align}
Its solution
\begin{align}
 \tfrac19\omega_2(q^3,p^3)
 ={}&\langle t|t\rangle
      +\langle q|p|t\rangle-\langle p|q|t\rangle
       \label{eq:vol2-cyclic-second-primitive}\\
    &-\langle q|t|p\rangle+\langle p|t|q\rangle
      +\langle t|q|p\rangle-\langle t|p|q\rangle\notag\\
    &-\langle q|q|p|p\rangle+\langle q|p|q|p\rangle
      +\langle q|p|p|q\rangle\notag\\
    &+\langle p|q|q|p\rangle-\langle p|q|p|q\rangle
      -\langle p|p|q|q\rangle\notag
\end{align}
also has zero Connes boundary and HKR projection.
Substitution in the raw differential verifies these identities:
internal terms replace $t$ by $qp-pq$, and adjacent mergers
cancel the cyclic faces. In particular the
$\langle t|t\rangle$ term is required in the second equation.
The contracted projection kernel extends these prescribed
lower components to all arities, even if a particular
row-reduction rule chooses different representatives.

\section{The finite matrix tangent pairing}
\label{sec:vol2-matrix-tangent-pairing}

The diagonal Koszul resolution also explains the nondegeneracy
of the orientation. Its top exterior product pairs the
resolution with its dual and gives
\[
 \operatorname{RHom}_{A^e}(A,A^e)\simeq A[-2].
\]
The volume $dq\wedge dp$ fixes this identification. Together
with its negative-cyclic lift, this is the smooth
$2$-Calabi--Yau orientation of the plane: the diagonal is a
perfect bimodule, and its bimodule dual is the displayed shift.
Extension and restriction along the quasi-isomorphism $R\to A$
transport this duality to $R$. Properness would require finite
total cohomology over $\mathbb C$, which $A$ does not have.
The orientation is a bimodule duality, distinct from a
nondegenerate scalar trace on a proper algebra.

Take the quotient by $GL_N$ of the derived affine scheme
represented by $R_N$. At a commuting pair $(X,Y)$ its tangent
complex is
\begin{equation}\label{eq:vol2-derived-matrix-tangent}
 \underbrace{\mathfrak{gl}_N}_{-1}
 \xrightarrow{\ d_0\ }
 \underbrace{\operatorname{Mat}_N\oplus\operatorname{Mat}_N}_{0}
 \xrightarrow{\ d_1\ }
 \underbrace{\mathfrak{gl}_N}_{1},
\end{equation}
with
\[
 d_0(a)=([a,X],[a,Y]),\qquad
 d_1(U,V)=[U,Y]+[X,V].
\]
The first arrow differentiates conjugation; the second
differentiates the relation. Their composite is
$[a,[X,Y]]=0$. Put
\begin{equation}\label{eq:vol2-matrix-tangent-form}
 \omega((U,V),(U',V'))=\operatorname{Tr}(UV'-VU')
\end{equation}
on the middle term. Pair the terms in degrees $-1$ and $1$
by $\operatorname{Tr}(ab)$, symmetrically under reversal as
required by graded skew-symmetry between two odd degrees.
Trace invariance gives
\[
 \omega(d_0a,(U,V))=\operatorname{Tr}(a\,d_1(U,V)).
\]
Thus these pairings form a degree-zero cochain pairing.
They are perfect term by term, including at points with
stabilizers. Discarding either outer term would discard part
of this self-dual tangent complex. This finite tangent
calculation and coherent cyclic invariance do not require
every chosen derivation homotopy to be Hamiltonian for one
fixed chain-level symplectic form.

\section{Deforming the coefficient relation}
\label{sec:vol2-coefficient-weyl-action}

Let $k_\varepsilon=\mathbb C[\varepsilon]$, with
$|\varepsilon|=0$, and define
\begin{equation}\label{eq:vol2-quantum-plane-resolution}
 \begin{gathered}
 W_\varepsilon=k_\varepsilon\langle q,p\rangle/
                              (qp-pq-\varepsilon),\\
 R_\varepsilon=k_\varepsilon\langle q,p,t\rangle,
 \qquad |t|=-1,\qquad dt=qp-pq-\varepsilon,
 \qquad dq=dp=0.
 \end{gathered}
\end{equation}
The projection $\pi_\varepsilon$ kills $t$.
Normal words $p^bq^a$ give a $k_\varepsilon$-basis of
$W_\varepsilon$: the replacement $qp\mapsto pq+\varepsilon$
terminates, and the leading pattern $qp$ has no overlap
ambiguity. They define a linear section $i_\varepsilon$.

Use the same marked-word contraction $h_0$ as before, with
$d_0t=qp$, but now put $\delta t=-pq-\varepsilon$.
Order expanded words first by length and then by the number
of pairs in which $q$ precedes $p$. Exchanging $qp$ and $pq$
decreases the second number; the constant term decreases the
first. Every term of $\delta h_0$ therefore decreases this
well-founded order.

\begin{proposition}\label{prop:vol2-integral-weyl-resolution}
The finite series
\[
 h_\varepsilon=h_0(1+\delta h_0)^{-1},\qquad
 p'=p_0(1+\delta h_0)^{-1}
\]
satisfy
\[
 dh_\varepsilon+h_\varepsilon d=1-i_\varepsilon\pi_\varepsilon,
 \quad h_\varepsilon^2=h_\varepsilon i_\varepsilon
                      =\pi_\varepsilon h_\varepsilon=0,
 \quad p'=\pi_\varepsilon.
\]
They extend continuously to the $\varepsilon$-adic completions,
allowing word support to grow with parameter order.
\end{proposition}
\begin{proof}
The marked-word proof of
Proposition~\ref{prop:vol2-plane-word-resolution} applies with
the decreasing order just described. Each surviving reduction
is precisely the PBW rule $qp\mapsto pq+\varepsilon$, so the
perturbed projection is $\pi_\varepsilon$. No division by
$\varepsilon$ occurs.

An element of the completed source is a series
$\sum_{r\geq0}\varepsilon^r a_r$ with each $a_r$ a finite
word polynomial. Every perturbation has nonnegative parameter
order. At a fixed output order, only finitely many earlier
coefficients contribute, and each has a terminating word
calculation. This defines the continuous extensions. The
identities hold modulo every power of $\varepsilon$, hence in
the separated inverse limit.
\end{proof}

Let $m(f\otimes g)=fg$ be ordinary polynomial multiplication.
On $k_\varepsilon[q,p]$, the Moyal product is
\begin{equation}\label{eq:vol2-coefficient-moyal-product}
 f\star_\varepsilon g
 =m\exp\!\left[\frac\varepsilon2
       (\partial_q\otimes\partial_p-\partial_p\otimes\partial_q)
       \right](f\otimes g).
\end{equation}
It is finite on polynomial inputs. On three tensor factors,
the three constant bidifferential operators commute; both
parenthesizations are the exponential of their sum. This
proves associativity. Symmetrization identifies this product
with multiplication in $W_\varepsilon$. For linear forms $u,v$,
centrality of their Weyl commutator gives
\[
 e^{\sigma u}e^{\sigma v}
   =e^{\varepsilon\{u,v\}/2}e^{\sigma(u+v)}.
\]
Differentiation in a formal scalar parameter proves this
exponential identity. Taking coefficients in $u,v$ gives the
polynomial product identity.

The commutator is divisible by $\varepsilon$ in the polynomial
ring. Write
\[
 B_\varepsilon(f,g)=
          (f\star_\varepsilon g-g\star_\varepsilon f)/\varepsilon,
 \qquad
 \mathfrak l_\varepsilon=k_\varepsilon[q,p]/k_\varepsilon1.
\]
Division here means the polynomial quotient; the parameter is
not inverted. The quotient is a Lie algebra. The cyclic
derivative formula \eqref{eq:vol2-hamiltonian-first-lift}
still defines a closed derivation of $R_\varepsilon$.
Under $\pi_\varepsilon$ it induces the strict Moyal
derivation $B_\varepsilon(f,-)$. Indeed its values on $q,p$
are the Weyl-ordered derivatives $-f_p,f_q$, and those
generators determine a derivation of the Weyl algebra.

\begin{theorem}\label{thm:vol2-coherent-coefficient-quantum-action}
There is an $L_\infty$ morphism over $k_\varepsilon$
\[
 F^\varepsilon:\mathfrak l_\varepsilon
       \longrightarrow\operatorname{Der}_{k_\varepsilon}
                                   (R_\varepsilon)
\]
with the specified first component. Each Taylor component
is polynomial and finite at fixed arity on polynomial inputs.
The action extends continuously to the parameter completions
and gives, by matrix substitution, a $GL_N$-equivariant action
on the derived representation algebra for every $N$.
\end{theorem}
\begin{proof}
Let $K_\varepsilon$ be the derivations whose generator values
lie in $\ker\pi_\varepsilon$. The differential of a
degree-$j$ derivation with values $(U,V,W)$ is still
\eqref{eq:vol2-derivation-differential}: derivations over
$k_\varepsilon$ kill the added scalar.
Replacing $h$ by $h_\varepsilon$ in
\eqref{eq:vol2-derivation-kernel-homotopy} therefore contracts
$K_\varepsilon$ in every degree.

Use the coalgebra signs
\eqref{eq:vol2-lie-coalgebra-signs} and replace the source
Poisson bracket by $B_\varepsilon$. The first action projects
to a strict action on $W_\varepsilon$. Every higher residual
is consequently closed and lies in $K_\varepsilon$.
Applying its contraction supplies the next component, exactly
as in Theorem~\ref{thm:vol2-all-arity-derived-action}.
Finitely many partitions and the terminating word order prove
polynomial finiteness. The maps never decrease parameter
order, giving the continuous completed identities. Matrix
substitution preserves differentials, brackets and
conjugation over the same base ring.
\end{proof}

The first parameter-dependent correction is already visible
on a cubic and a quartic:
\begin{equation}\label{eq:vol2-first-quantum-hamiltonian-contact}
 \begin{aligned}
 B_\varepsilon(q^3,p^4)&=12q^2p^3+6\varepsilon^2p,\\
 F^\varepsilon_2(q^3,p^4)q
     &=6(-pqt+ptq-qtp+tpq-\varepsilon t),\\
 F^\varepsilon_2(q^3,p^4)p&=6(-p^2t+tp^2),\\
 F^\varepsilon_2(q^3,p^4)t&=6(pt^2+tpt).
 \end{aligned}
\end{equation}
The first and third odd Moyal terms give the bracket in the
first line. Applying $\partial=[d,-]$ to the remaining lines
gives the binary defect on all three generators. In particular
$d(-6\varepsilon t)$ contributes $+6\varepsilon^2$ to the
$q$-component, matching the contribution of
$-F_1(6\varepsilon^2p)$ to the defect.
The third component on the ordered triple
$(q^3,p^4,q^2p)$ is
\begin{equation}\label{eq:vol2-quantum-ternary-action}
 F^\varepsilon_3q=-12tqt+12t^2q,
 \qquad F^\varepsilon_3p=-12pt^2,
 \qquad F^\varepsilon_3t=-6t^3.
\end{equation}
Substitution of $dt=qp-pq-\varepsilon$ verifies its ternary
equation with the binary terms in
\eqref{eq:vol2-first-quantum-hamiltonian-contact}.

The classical positive-weight completion and this parameter
completion use different estimates. Higher Moyal terms can
lower polynomial degree, as the term $\varepsilon^2p$ shows.
The parameter-adic theorem does not assert that the classical
pointed polynomial subspace remains closed under the deformed
bracket.

There is a relative negative-cyclic orientation over the
parameter ring. Form normalized chains over $k_\varepsilon$,
so that a parameter multiple of $1$ is zero in a bar slot.
In the raw totalization,
\begin{equation}\label{eq:vol2-relative-quantum-orientation}
 \omega_\varepsilon=1[q|p]-1[p|q]-1[t]
\end{equation}
satisfies
\[
 d_{\mathrm H}\omega_\varepsilon=-\varepsilon1[1]=0,
 \qquad B\omega_\varepsilon=0.
\]
This is a relative mixed-complex assertion. It uses the
parameter in the base ring; an absolute de Rham variation
of that parameter is additional data. The cycle itself also
does not supply an equivariant cyclic comparison to a
matrix-quantized recipient.

\section{Matrix quantization and finite representations}
\label{sec:vol2-two-quantizations}

Matrix quantization starts with the canonical matrix coordinates,
using an independent parameter $\hbar$. Its even Weyl algebra
$W_{N,\hbar}$ over $\mathbb C[\hbar]$ has
\[
 [X_{ij},Y_{kl}]=\hbar\delta_{il}\delta_{jk},\qquad
 [X_{ij},X_{kl}]=[Y_{ij},Y_{kl}]=0.
\]
Weyl ordering is symmetrization of all coordinate factors.
The Moyal formula on this finite-dimensional symplectic vector
space identifies its multiplication with the corresponding
constant-tensor star product. The commutator of two scalar
quadratic polynomials, divided by $\hbar$, is exactly their
Poisson bracket: the third and higher odd Moyal terms
differentiate each quadratic at least three times and vanish.
Thus the strict scalar quadratic sector constructed above
remains strict under matrix quantization.

The ordering of the matrix moment is part of this assertion.
As operator-valued matrices, set
\begin{equation}\label{eq:vol2-ordered-matrix-moment}
 \widehat\mu=XY-YX-N\hbar I_N.
\end{equation}
The Weyl symbols of $XY$ and $YX$ differ from their commuting
symbols by $N\hbar I_N/2$ and $-N\hbar I_N/2$, respectively.
Consequently $\widehat\mu$ has Weyl symbol $[X,Y]$ in
commuting coordinates. In particular
\[
 \operatorname{Tr}(XY-YX)=N^2\hbar,
 \qquad \operatorname{Tr}\widehat\mu=0.
\]
This scalar ordering correction is fixed by the coordinate
commutators. A specified nonzero value of a moment map is a
separate datum.

The full quantum BRST algebra also retains both ghost degrees.
Choose a basis $a_i$ of $\mathfrak{gl}_N$ with
$[a_i,a_j]=\sum_k f_{ij}^ka_k$, and set
$\ell_i=\operatorname{Tr}(a_i\widehat\mu)$.
These are the Weyl-ordered quadratic comoments, with
\[
 [\ell_i,\ell_j]=\hbar\sum_k f_{ij}^k\ell_k.
\]
This follows either by canonical matrix differentiation and
exact quadratic quantization, or by realizing
$X_{ij}=x_{ij}$ and $Y_{ij}=-\hbar\partial_{x_{ji}}$:
$\ell_a=\hbar\tau(a)$, where
$\tau(a)X=[X,a]$ is the infinitesimal pullback action.

Adjoin odd generators $c^i,b_i$ of degrees $1,-1$, commuting
with the even matter generators and satisfying
\[
 c^ic^j+c^jc^i=0,\qquad b_ib_j+b_jb_i=0,\qquad
 b_ic^j+c^jb_i=\hbar\delta_i^j.
\]
Write $\mathcal B_{N,\hbar}$ for the resulting graded tensor
product with $W_{N,\hbar}$. Define
\begin{equation}\label{eq:vol2-quantum-bfv-hamiltonian}
 \Omega=\sum_i c^i\ell_i
          -\frac12\sum_{i,j,k}f_{ij}^kc^ic^jb_k,
 \qquad Q_\hbar=\hbar^{-1}[\Omega,-]_{\mathrm{gr}}.
\end{equation}
Modulo $\hbar$ the algebra is supercommutative, so the
commutator is divisible by $\hbar$ and the formula is integral.
Under the trace identification of the ghost and its dual,
its classical limit is
\eqref{eq:vol2-classical-bfv-differential}.

\begin{proposition}\label{prop:vol2-integral-quantum-brst}
The element $\Omega$ satisfies $\Omega^2=0$.
Consequently $Q_\hbar$ is a degree-one derivation with
$Q_\hbar^2=0$ over $\mathbb C[\hbar]$.
\end{proposition}
\begin{proof}
Normal words, with all $c$ generators before all $b$ generators,
span the ghost algebra. Represent $c^i$ by exterior
multiplication and $b_i$ by $\hbar\partial_{c^i}$, with left
odd derivatives. After inverting $\hbar$, their normal words
form a basis of the endomorphisms of the exterior space.
They are therefore linearly independent over
$\mathbb C[\hbar]$ as well. Together with the Weyl basis this
proves parameter-torsion freeness and faithfulness on the
tensor product of the matter regular module and the ghost
exterior module.

In this representation,
\[
 \Omega=\sum_i c^i\ell_i+\hbar d_{\mathrm{CE}},\qquad
 d_{\mathrm{CE}}c^k=-\frac12\sum_{i,j}f_{ij}^kc^ic^j.
\]
Jacobi gives $d_{\mathrm{CE}}^2=0$. The square of the first
term is
$\frac{\hbar}{2}\sum f_{ij}^kc^ic^j\ell_k$; its mixed term
with $\hbar d_{\mathrm{CE}}$ is the negative of this expression.
Thus $\Omega^2=0$ in a faithful representation and in the
algebra. For odd $\Omega$, the double graded commutator with
$\Omega$ is the commutator with $\Omega^2$.
Hence $\hbar^2Q_\hbar^2=0$, and torsion freeness gives the result.
\end{proof}

The full gauge complex has extra states even at rank one.
Scalar conjugation is trivial, so $\ell=0$, the Lie bracket
vanishes and $Q_\hbar=0$. In its degree-zero ghost algebra,
$1$ and $u=cb$ are independent and obey
\[
 u^2=\hbar u.
\]
Thus the full BRST cohomology is not obtained by simply
replacing it with the ordinary one-pair Weyl reduction.
Removing the ineffective gauge direction changes the input.

In contrast, a finite-dimensional representation of the
\emph{coefficient} Weyl algebra has commuting coordinate
entries and the relation $[X,Y]=\varepsilon I_N$.
Its derived representation algebra is
\begin{equation}\label{eq:vol2-coefficient-derived-representations}
 R_{N,\varepsilon}
   =\mathbb C[X_{ij},Y_{ij},\varepsilon]\otimes\Lambda(\Psi_{ij}),
 \qquad Q\Psi=[X,Y]-\varepsilon I_N,
 \qquad QX=QY=Q\varepsilon=0.
\end{equation}
This is the recipient supplied by matrix substitution in
Theorem~\ref{thm:vol2-coherent-coefficient-quantum-action}.

\begin{proposition}\label{prop:vol2-finite-weyl-representations}
For $N>0$, multiplication by $\varepsilon$ is null-homotopic
on $R_{N,\varepsilon}$. Its cohomology is annihilated by
$\varepsilon$, and its localization at $\varepsilon$ is
contractible.
\end{proposition}
\begin{proof}
Commuting entries give $\operatorname{Tr}[X,Y]=0$.
The element $v=-\operatorname{Tr}\Psi/N$ has degree $-1$
and satisfies $Qv=\varepsilon$. Multiplication by $v$ obeys
\[
 Q(va)+vQa=\varepsilon a
\]
for every homogeneous $a$. After inverting $\varepsilon$,
multiplication by $v/\varepsilon$ contracts the complex.
\end{proof}

This contraction is integral before localization. It does not
contract the special fibre. At rank one the complex is
\[
 \mathbb C[x,y,\varepsilon]\psi
     \xrightarrow{-\varepsilon}
 \mathbb C[x,y,\varepsilon],\qquad |\psi|=-1.
\]
Its cohomology is $\mathbb C[x,y]$ in degree zero, annihilated
by $\varepsilon$. Derived specialization at $\varepsilon=0$
is computed by the displayed free complex and retains the
degree-minus-one class $\psi$. Simply specializing its
degree-zero cohomology would miss that class.

The equations \eqref{eq:vol2-ordered-matrix-moment} and
\eqref{eq:vol2-coefficient-derived-representations} therefore
cannot be interchanged by naming their parameters alike.
One quantizes the matrix coordinate algebra; the other asks
for matrix representations of a quantized coefficient algebra.
A comparison between them must specify a map and its effect
on relation classes.

\section{Framed and infinite-dimensional recipients}
\label{sec:vol2-weyl-action-recipients}

Framing changes the finite-dimensional equation. Add a column
$\mathbf i\in\mathbb C^N$ and a row
$\mathbf j\in(\mathbb C^N)^*$, and use
\begin{equation}\label{eq:vol2-framed-moment-equation}
 [X,Y]+\mathbf i\mathbf j=\varepsilon I_N.
\end{equation}
Simultaneous conjugation extends by
$\mathbf i\mapsto g\mathbf i$ and
$\mathbf j\mapsto\mathbf jg^{-1}$. The extra rank-one
term permits a nonzero trace of the prescribed moment value.

\begin{proposition}\label{prop:vol2-framed-weyl-chart}
On the ordered chart with pairwise distinct $x_i$, the formulas
\begin{equation}\label{eq:vol2-framed-weyl-coordinates}
 \begin{gathered}
 X=\operatorname{diag}(x_i),\qquad
 \mathbf i_i=1,\qquad\mathbf j_i=\varepsilon,\\
 Y_{ii}=y_i,\qquad
 Y_{ij}=-\frac{\varepsilon}{x_i-x_j}\quad(i\ne j)
 \end{gathered}
\end{equation}
give a free framed orbit and a smooth local quotient of
relative dimension $2N$ over the parameter. At fixed
$\varepsilon$, the reduced symplectic form is
$\sum_i dx_i\wedge dy_i$ on this ordered cover.
\end{proposition}
\begin{proof}
The off-diagonal entries of the left side of
\eqref{eq:vol2-framed-moment-equation} are
$-\varepsilon+\varepsilon=0$; its diagonal entries are
$\varepsilon$. A stabilizer commuting with $X$ is diagonal,
and fixing the column with all entries one makes it the
identity. The infinitesimal action is injective.

On the open set with simple $X$ and nonzero components of
$\mathbf i$ in its eigenbasis, diagonalize $X$ and use diagonal
gauge transformations to set those components equal to one.
The moment equation then forces the displayed values of
$\mathbf j$ and the off-diagonal entries of $Y$.
Thus $(x_i,y_i)$ actually parametrize the local quotient.

With the canonical form
$\operatorname{Tr}(dX\wedge dY)
 +\sum_i d\mathbf i_i\wedge d\mathbf j_i$,
the moment differential is the symplectic transpose of the
infinitesimal action, so it is surjective. The quotient is
therefore smooth of relative dimension
$(2N^2+2N)-2N^2=2N$ near this orbit.
On the slice, $d\mathbf i=d\mathbf j=0$ at fixed parameter,
and $dX$ is diagonal. The pullback of the form is exactly
$\sum_i dx_i\wedge dy_i$. Permuting the eigenvalues gives
the remaining finite identification.
\end{proof}

An infinite-dimensional unframed recipient keeps the original
coefficient equation. Complete $W_\varepsilon$ in the
$\varepsilon$-adic topology and write it as
$\widehat W_\varepsilon$. Use complete left modules with
continuous maps and strict exact sequences, resolved by
topologically free modules. The completed polynomial Fock module
\[
 \mathcal P=\mathbb C[p][[\varepsilon]],\qquad
 q\cdot f=\varepsilon\partial_pf,\qquad p\cdot f=pf
\]
satisfies $[q,p]=\varepsilon$. A coefficient at each parameter
order is a polynomial in $p$, with no uniform degree bound.

\begin{proposition}\label{prop:vol2-integral-fock-self-extension}
Continuous derived Hom in this exact category gives
\[
 \operatorname{Ext}^0_{\widehat W_\varepsilon}(\mathcal P,\mathcal P)
     =\mathbb C[[\varepsilon]],
 \qquad
 \operatorname{Ext}^1_{\widehat W_\varepsilon}(\mathcal P,\mathcal P)
     =\mathbb C[p],
\]
with all other degrees zero. Multiplication by $\varepsilon$
annihilates the degree-one group.
\end{proposition}
\begin{proof}
In the completed PBW basis $p^bq^a$, right multiplication by
$q$ increases $a$ by one. It is injective with closed image,
and its cokernel is $\mathcal P$. As a map of left modules
it gives the strict topologically free resolution
\[
 0\longrightarrow\widehat W_\varepsilon
   \xrightarrow{R_q}\widehat W_\varepsilon
   \longrightarrow\mathcal P\longrightarrow0.
\]
Continuous Hom into $\mathcal P$ is the two-term complex
\[
 \mathcal P\xrightarrow{\varepsilon\partial_p}\mathcal P
\]
in degrees zero and one. Coefficientwise polynomial
integration is a continuous right inverse to $\partial_p$.
The kernel is $\mathbb C[[\varepsilon]]$, the image is
$\varepsilon\mathcal P$, and the cokernel is
$\mathcal P/\varepsilon\mathcal P=\mathbb C[p]$.
\end{proof}

The degree-one self-extension disappears when the parameter
is inverted. Its integral presence records information that
generic specialization loses. The framed chart and the Fock
module thus avoid the finite unframed contraction in different
ways: the first changes the equation, and the second changes
the dimension and completion of the module.

\section{Hamiltonian and geometric comparison maps}
\label{sec:vol2-derived-hamiltonian-comparison}

The construction has produced an action on an associative
resolution, an action on derived matrix representations,
and a coherent classical negative-cyclic orientation.
The quotient-by-constants obstruction fixes the source of
the stronger Hamiltonian problem. Retain either the full
classical Poisson algebra $A$ with $J_N(1)=N$, or its pointed
Lie subalgebra, and keep the scalar quadratic map fixed.

\begin{problem}\label{prob:vol2-derived-hamiltonian-moment-map}
Construct a homotopy moment map to the full classical BFV
algebra whose first component is $J_N$, whose cubic binary
component is $-\tfrac32\operatorname{Tr}(\Psi[X,Y])$ up to
specified homotopies, and whose choices respect block-diagonal
sums of matrices. Compare its Hamiltonian derivation action
with Theorem~\ref{thm:vol2-all-arity-derived-action}, retaining
the cyclic orientation.

For matrix quantization, construct a filtered derived
comparison on the full quantum BRST complex that agrees
with the ordinary quantum trace action in degree zero and
preserves the relation classes, stabilizer states and cyclic
operations. Require the prescribed classical limit and cubic
primitive, and extend the comparison across the matrix
discriminant. Determine whether the coefficient-Weyl action
compares to this recipient, to a framed recipient, or to an
infinite-dimensional module by maps preserving those data.
\end{problem}

The positive gauge ghosts in
\eqref{eq:vol2-classical-bfv-differential} already belong to
the action constructed here. Supplying them is not the
missing Hamiltonian comparison. Likewise a contraction of
one finite coefficient representation does not obstruct
every framed, completed or quantum-reduced recipient.

For mixed holomorphic--topological theory there is a further
geometric map to construct. The polynomial cyclic operations
must be realized by supported operations of the independently
defined field theory. Such a comparison must preserve the
boundary polarization, parameter normalization, descendant
operations and sewing. At four inputs it must include the
simultaneous collision support
$\Delta_{13}\cap\Delta_{24}$ from
Chapter~\ref{ch:vol2-supported-compositions}.
The universal word contraction solves the algebraic action
equations; the geometric comparison must carry their
solutions to these specified supports and operations.
